\begin{table}[h]
\centering
\caption{Summary Statistics by Zoning Designation}
\label{tab:summary_stats_zone}
\begin{threeparttable}
\begin{tabular*}{1.0\textwidth}{@{\extracolsep{\fill}}l*{2}{r}}
\toprule
 & Industrial & Residential \\
\midrule
Residents & \makecell[tr]{55.85  (61.49)} & \makecell[tr]{51.44  (47.40)} \\
Households & \makecell[tr]{16.22  (20.42)} & \makecell[tr]{14.35  (13.85)} \\
Median Rent & \makecell[tr]{28.28  (172.09)} & \makecell[tr]{34.37  (182.06)} \\
Median Home Value & \makecell[tr]{1943.88  (2313.48)} & \makecell[tr]{3289.95  (2835.71)} \\
Percent Black & \makecell[tr]{0.26  (0.39)} & \makecell[tr]{0.14  (0.32)} \\
Share of Black Residents & \makecell[tr]{0.00  (0.00)} & \makecell[tr]{0.00  (0.00)} \\
Highway Present (1940) & \makecell[tr]{0.14  (0.35)} & \makecell[tr]{0.07  (0.25)} \\
Highway Present (1959) & \makecell[tr]{0.16  (0.37)} & \makecell[tr]{0.08  (0.27)} \\
Highway Constructed (1940-1959) & \makecell[tr]{0.02  (0.15)} & \makecell[tr]{0.01  (0.11)} \\
Total Households & 34129 & 67645 \\
Total Squares & 2104 & 4713 \\
\bottomrule
\end{tabular*}
\begin{tablenotes}[flushleft]
\footnotesize
\item This table contains summary statistics for the full sample of grid squares, split by zoning designation. The table reports the mean and standard deviation (in parentheses) for each variable, as well as the total number of households in each zoning category.
\end{tablenotes}
\end{threeparttable}
\end{table}
